% GENERATED FILE. Edit canonical lessons, then run npm run generate:books.
% canonical-source-hash: fnv1a64:70823dbaa0b62e74
% canonical-lessons: FR-C473-fierte, FR-C473-orgueil, FR-C473-vaniteux, FR-C473-modeste, FR-C473-vanter, FR-R473-revision-466-473

\chapter{Fierté et modestie}
\label{ch:fr-fierte-et-modestie}
\hlchaptermodality{473}

\begin{chapteropening}
\textbf{By the end of this chapter:} \emph{I can talk about pride, vanity and modesty.}

Complete the last lesson of chapter 473: I can talk about pride, vanity and modesty.
\end{chapteropening}

\section[la fierté]{la fierté — pride}

\label{lesson:FR-C473-fierte}



\begin{quote}
Before the new one: say the French for to move, then the French for to deeply upset.
\end{quote}

\subsection*{You'll want to know: la fierté}
\textbf{la fierté} — \textquotedblleft{}pride\textquotedblright{}.

\begin{quote}
Mon père parle de mon travail avec fierté.
\end{quote}

\begin{quote}
\emph{My father speaks about my work with pride.}
\end{quote}

Built on \emph{fier} \textquotedblleft{}proud\textquotedblright{}; a good kind of pride.

\begin{cousinweb}
\emph{la fierté} is built on \emph{fier} \textquotedblleft{}proud\textquotedblright{} with \emph{-té} (EtymDB). \emph{Fier} is inherited from Latin \emph{ferus} \textquotedblleft{}wild, untamed\textquotedblright{} (Lewis \& Short).
\end{cousinweb}

\begin{grammarlens}[title={What you've built}]
One more word for telling, explaining and describing at B1.
\end{grammarlens}

\subsection*{Your turn}
\begin{itemize}
\raggedright
  \item \emph{Say it:} \emph{la fierté}
  \item \emph{Say it:} \emph{la fierté}, once more
  \item \emph{Say it:} a sentence of your own with \emph{la fierté}
  \item \emph{Recall:} say the French for to move, then the French for to deeply upset, then say \emph{la fierté} again
\end{itemize}

\begin{tcolorbox}[breakable,colback=teal!4,colframe=teal!35!black,title={Before you move on}]
What does la fierté mean? (\textquotedblleft{}Pride\textquotedblright{}.) Say it once more.
\end{tcolorbox}
\section[l'orgueil]{l'orgueil — arrogance, excessive pride}

\label{lesson:FR-C473-orgueil}



\begin{quote}
Before the new one: say the French for to deeply upset, then the French for pride.
\end{quote}

\subsection*{You'll want to know: l'orgueil}
\textbf{l'orgueil} — \textquotedblleft{}arrogance, excessive pride\textquotedblright{}.

\begin{quote}
Par orgueil, il n'a pas demandé d'aide.
\end{quote}

\begin{quote}
\emph{Out of pride, he did not ask for help.}
\end{quote}

Masculine: \emph{un orgueil}. Stronger and more negative than \emph{la fierté}.

\begin{cousinweb}
\emph{l'orgueil} comes from Frankish \emph{urgōli} \textquotedblleft{}pride\textquotedblright{}, a Germanic word related to Old English \emph{orgel} \textquotedblleft{}pride\textquotedblright{} (EtymDB).
\end{cousinweb}

\begin{grammarlens}[title={What you've built}]
One more word for telling, explaining and describing at B1.
\end{grammarlens}

\subsection*{Your turn}
\begin{itemize}
\raggedright
  \item \emph{Say it:} \emph{l'orgueil}
  \item \emph{Say it:} \emph{l'orgueil}, once more
  \item \emph{Say it:} a sentence of your own with \emph{l'orgueil}
  \item \emph{Recall:} say the French for to deeply upset, then the French for pride, then say \emph{l'orgueil} again
\end{itemize}

\begin{tcolorbox}[breakable,colback=teal!4,colframe=teal!35!black,title={Before you move on}]
What does l'orgueil mean? (\textquotedblleft{}Arrogance, excessive pride\textquotedblright{}.) Say it once more.
\end{tcolorbox}
\section[vaniteux]{vaniteux — vain, conceited}

\label{lesson:FR-C473-vaniteux}



\begin{quote}
Before the new one: say the French for pride, then the French for arrogance.
\end{quote}

\subsection*{You'll want to know: vaniteux}
\textbf{vaniteux} — \textquotedblleft{}vain, conceited\textquotedblright{}.

\begin{quote}
Il est vaniteux et parle toujours de lui.
\end{quote}

\begin{quote}
\emph{He is vain and always talks about himself.}
\end{quote}

Feminine \emph{vaniteuse}. The noun is \emph{la vanité}.

\begin{cousinweb}
\emph{vaniteux} is built on \emph{vanité}, a learned borrowing from Latin \emph{vānitās} \textquotedblleft{}emptiness, vanity\textquotedblright{}, from \emph{vānus} \textquotedblleft{}empty\textquotedblright{} (Lewis \& Short).
\end{cousinweb}

\begin{grammarlens}[title={What you've built}]
One more word for telling, explaining and describing at B1.
\end{grammarlens}

\subsection*{Your turn}
\begin{itemize}
\raggedright
  \item \emph{Say it:} \emph{vaniteux}
  \item \emph{Say it:} \emph{vaniteux}, once more
  \item \emph{Say it:} a sentence of your own with \emph{vaniteux}
  \item \emph{Recall:} say the French for pride, then the French for arrogance, then say \emph{vaniteux} again
\end{itemize}

\begin{tcolorbox}[breakable,colback=teal!4,colframe=teal!35!black,title={Before you move on}]
What does vaniteux mean? (\textquotedblleft{}Vain, conceited\textquotedblright{}.) Say it once more.
\end{tcolorbox}
\section[modeste]{modeste — modest}

\label{lesson:FR-C473-modeste}



\begin{quote}
Before the new one: say the French for arrogance, then the French for vain.
\end{quote}

\subsection*{You'll want to know: modeste}
\textbf{modeste} — \textquotedblleft{}modest\textquotedblright{}.

\begin{quote}
Elle a gagné le prix, mais elle reste modeste.
\end{quote}

\begin{quote}
\emph{She won the prize, but she remains modest.}
\end{quote}

Also simple, small: \emph{un salaire modeste}.

\begin{cousinweb}
\emph{modeste} is a learned borrowing from Latin \emph{modestus} \textquotedblleft{}keeping within measure\textquotedblright{}, from \emph{modus} \textquotedblleft{}measure\textquotedblright{} (Lewis \& Short).
\end{cousinweb}

\begin{grammarlens}[title={What you've built}]
One more word for telling, explaining and describing at B1.
\end{grammarlens}

\subsection*{Your turn}
\begin{itemize}
\raggedright
  \item \emph{Say it:} \emph{modeste}
  \item \emph{Say it:} \emph{modeste}, once more
  \item \emph{Say it:} a sentence of your own with \emph{modeste}
  \item \emph{Recall:} say the French for arrogance, then the French for vain, then say \emph{modeste} again
\end{itemize}

\begin{tcolorbox}[breakable,colback=teal!4,colframe=teal!35!black,title={Before you move on}]
What does modeste mean? (\textquotedblleft{}Modest\textquotedblright{}.) Say it once more.
\end{tcolorbox}
\section[vanter]{vanter — to praise; (se vanter) to boast}

\label{lesson:FR-C473-vanter}



\begin{quote}
Before the new one: say the French for vain, then the French for modest.
\end{quote}

\subsection*{You'll want to know: vanter}
\textbf{vanter} — \textquotedblleft{}to praise; (se vanter) to boast\textquotedblright{}.

\begin{quote}
Mon collègue se vante de sa nouvelle voiture.
\end{quote}

\begin{quote}
\emph{My colleague is boasting about his new car.}
\end{quote}

Mostly reflexive: \emph{se vanter de quelque chose}. \emph{Vanter un produit} — to praise a product.

\begin{cousinweb}
\emph{vanter} is inherited from Old French \emph{vanter} (EtymDB), usually traced to a Late Latin \emph{vānitāre} \textquotedblleft{}to speak emptily\textquotedblright{}, built on Latin \emph{vānus} \textquotedblleft{}empty, vain\textquotedblright{} (Lewis \& Short).
\end{cousinweb}

\begin{grammarlens}[title={What you've built}]
One more word for telling, explaining and describing at B1.
\end{grammarlens}

\subsection*{Your turn}
\begin{itemize}
\raggedright
  \item \emph{Say it:} \emph{vanter}
  \item \emph{Say it:} \emph{vanter}, once more
  \item \emph{Say it:} a sentence of your own with \emph{vanter}
  \item \emph{Recall:} say the French for vain, then the French for modest, then say \emph{vanter} again
\end{itemize}

\begin{tcolorbox}[breakable,colback=teal!4,colframe=teal!35!black,title={Before you move on}]
What does vanter mean? (\textquotedblleft{}To praise; (se vanter) to boast\textquotedblright{}.) Say it once more.
\end{tcolorbox}
\section[(dialogue)]{Sentiments et santé}

\label{lesson:FR-R473-revision-466-473}



\begin{quote}
Say the French for shyness and sociable.
\end{quote}

\subsection*{Your turn}
Say these aloud:

\begin{itemize}
\raggedright
  \item shyness, sociable, reserved, to sulk, to tease
  \item stubborn, laziness, stingy, clumsy, scatterbrained
  \item sincere, liar, naive, cunning, to lie
  \item mood, calm, moody, to annoy, to irritate
  \item fright, to frighten, to jump, relief, to reassure
  \item disappointment, to disappoint, sorrow, to sob, to comfort
  \item nostalgia, to regret, melancholy, to move, to deeply upset
  \item pride, arrogance, vain, modest, to praise
\end{itemize}

\begin{tcolorbox}[breakable,colback=teal!4,colframe=teal!35!black,title={Before you move on}]
Shyness? (\textbf{la timidité}.) Sociable? (\textbf{sociable}.) Reserved? (\textbf{réservé}.) To sulk? (\textbf{bouder}.) To tease? (\textbf{taquiner}.) Stubborn? (\textbf{têtu}.) Laziness? (\textbf{la paresse}.) Stingy? (\textbf{avare}.) Clumsy? (\textbf{maladroit}.) Scatterbrained? (\textbf{étourdi}.) Sincere? (\textbf{sincère}.) Liar? (\textbf{le menteur}.) Naive? (\textbf{naïf}.) Cunning? (\textbf{rusé}.) To lie? (\textbf{mentir}.) Mood? (\textbf{l'humeur}.) Calm? (\textbf{serein}.) Moody? (\textbf{lunatique}.) To annoy? (\textbf{énerver}.) To irritate? (\textbf{agacer}.) Fright? (\textbf{la frayeur}.) To frighten? (\textbf{effrayer}.) To jump? (\textbf{sursauter}.) Relief? (\textbf{le soulagement}.) To reassure? (\textbf{rassurer}.) Disappointment? (\textbf{la déception}.) To disappoint? (\textbf{décevoir}.) Sorrow? (\textbf{le chagrin}.) To sob? (\textbf{sangloter}.) To comfort? (\textbf{consoler}.) Nostalgia? (\textbf{la nostalgie}.) To regret? (\textbf{regretter}.) Melancholy? (\textbf{la mélancolie}.) To move? (\textbf{émouvoir}.) To deeply upset? (\textbf{bouleverser}.) Pride? (\textbf{la fierté}.) Arrogance? (\textbf{l'orgueil}.) Vain? (\textbf{vaniteux}.) Modest? (\textbf{modeste}.) To praise? (\textbf{vanter}.)
\end{tcolorbox}
